% Apéndices

\section{Anexo A: Instalación y configuración del entorno}

\begin{lstlisting}[language=bash, caption={Instalación del entorno de desarrollo.}]
# Clonar el repositorio
git clone <repo_url>
cd qsimov

# Crear entorno virtual con Python 3.8
python3.8 -m venv venv
source venv/bin/activate

# Instalar dependencias
pip install -r requirements.txt

# Compilar extensiones Cython (obligatorio antes del primer uso)
python setup.py build_ext --inplace
\end{lstlisting}

\section{Anexo B: Ejemplo de uso completo (Keras, aprendizaje continual)}

\begin{lstlisting}[language=Python, caption={Ejemplo completo de aprendizaje continual con Keras.}]
import numpy as np
import tensorflow.keras as kr
from qsimov.keras_path_selector import KerasPathSelector
from qsimov.keras_qsimov_linear_system import KerasQsimovLinearSystem

# 1. Pre-entrenar modelo con datos de la fase 1
model = kr.Sequential([
    kr.layers.Conv2D(32, 3, activation='relu', input_shape=(28, 28, 1)),
    kr.layers.MaxPooling2D(),
    kr.layers.Flatten(),
    kr.layers.Dense(16, activation='relu'),
    kr.layers.Dense(10, activation='softmax'),
])
model.compile(optimizer='adam',
              loss='sparse_categorical_crossentropy',
              metrics=['accuracy'])
model.fit(x_phase1, y_phase1, epochs=10)

# 2. Crear PathSelector sobre el modelo entrenado
path_selector = KerasPathSelector(
    neural_network=model,
    initial_layer=-2,   # dos ultimas capas Dense
    verbose=1
)

# 3. Entrenar QsimovLinearSystem con fase 1
qls = KerasQsimovLinearSystem(
    path_selector=path_selector,
    solver="back_substitution",
    absolute_cutoff=1e-2,
    relative_cutoff=1e6,
)
qls.fit(x_phase1, y_phase1, batch_size=256)

# 4. Actualizar con fase 2 sin olvido (NO llamar reset_equations())
qls.fit(x_phase2, y_phase2, batch_size=256)

# 5. Evaluar en ambas fases
acc_old = np.mean(qls.predict(x_test_phase1).argmax(1) == y_test_phase1)
acc_new = np.mean(qls.predict(x_test_phase2).argmax(1) == y_test_phase2)
print(f"Precision clases antiguas: {acc_old:.3f}")
print(f"Precision clases nuevas:   {acc_new:.3f}")
\end{lstlisting}

\section{Anexo C: Guía de reproducibilidad}

Esta sección describe de forma precisa el entorno, la secuencia de instalación y los comandos
exactos necesarios para reproducir todos los experimentos presentados en este documento.

\subsection{C.1 Entorno de ejecución}

\textbf{Hardware}: Servidor \texttt{fermidi1-informat-uv-es}, Universitat de València. GPU NVIDIA
disponible en el nodo de cómputo.

\textbf{Software}:

\begin{table}[htbp]
  \centering
  \caption{Versiones del entorno software.}
  \label{tab:software}
  \begin{tabular}{ll}
    \toprule
    \textbf{Componente} & \textbf{Versión} \\
    \midrule
    Python              & 3.10 \\
    PyTorch             & 2.x (con CUDA) \\
    TensorFlow / Keras  & 2.11 (solo para experimentos Keras) \\
    NumPy               & 1.21 \\
    Cython              & 3.0 \\
    MLflow              & 2.x \\
    Plotly              & 5.x \\
    \bottomrule
  \end{tabular}
\end{table}

\subsection{C.2 Instalación del entorno}

\begin{lstlisting}[language=bash, caption={Configuración del entorno desde cero.}]
# 1. Clonar el repositorio del proyecto
git clone <url_repositorio>
cd qsimov

# 2. Crear entorno virtual Python 3.10
python3.10 -m venv venv
source venv/bin/activate

# 3. Instalar dependencias
pip install -r requirements.txt

# 4. Compilar las extensiones Cython (OBLIGATORIO antes del primer uso)
#    Sin este paso, las importaciones del modulo paths/ fallaran.
python setup.py build_ext --inplace
\end{lstlisting}

\subsection{C.3 Descarga y preprocesado de datos}

Cada experimento descarga y preprocesa sus propios datos automáticamente si no se especifica
\texttt{--skip-data-preparation}. Los datasets se almacenan en \texttt{data/} dentro del directorio
del proyecto. Para ImageNet es necesario proporcionar los archivos descargados previamente (ver
\texttt{experiments/imagenet\_subset\_by\_splits/preprocess\_data.py}).

\subsection{C.4 Comandos de ejecución por experimento}

Todos los comandos se ejecutan desde la raíz del directorio \texttt{qsimov/} con el entorno virtual
activo. Los resultados (modelos, historiales, gráficas HTML) se guardan en \texttt{results/} y un
resumen del experimento en \texttt{experiments/<nombre>/run\_exports/}.

\textbf{Experimento 1: Velocidad y precisión (MNIST)}

\begin{lstlisting}[language=bash]
python3 experiments/mnist_speed_loss/main.py \
    --framework pytorch \
    --processor gpu \
    --epochs 10
\end{lstlisting}

Resultados en: \path{results/mnist/learning_rate_speed_loss/pytorch/gpu/}

\textbf{Experimento 2: Robustez al learning rate (MNIST)}

\begin{lstlisting}[language=bash]
python3 experiments/mnist_learning_rate/main.py \
    --framework pytorch \
    --processor gpu \
    --epochs 10
\end{lstlisting}

Resultados en: \path{results/mnist/learning_rate_speed_loss/pytorch/gpu/}

\textbf{Experimento 3: Olvido catastrófico (MNIST)}

\begin{lstlisting}[language=bash]
python3 experiments/mnist_forgetting/main.py \
    --framework pytorch \
    --processor gpu \
    --base-epochs 10 \
    --epochs 5
\end{lstlisting}

Resultados en: \texttt{results/mnist/forgetting/pytorch/gpu/}

\textbf{Experimento 4: Datos escasos y splits (CIFAR-10)}

\begin{lstlisting}[language=bash]
python3 experiments/cifar10_gradient_by_splits/main.py \
    --framework pytorch \
    --processor gpu \
    --model_name lenet \
    --initial-layer -1 \
    --epochs 10 \
    --splits 100 1000 5000
\end{lstlisting}

Resultados en: \texttt{results/cifar10/gradient\_by\_splits/pytorch/gpu/}

\textbf{Experimento 5: Olvido catastrófico (CIFAR-10)}

\begin{lstlisting}[language=bash]
python3 experiments/cifar10_forgetting/main.py \
    --framework pytorch \
    --processor gpu \
    --base-epochs 20 \
    --epochs 5
\end{lstlisting}

Resultados en: \texttt{results/cifar10/forgetting/pytorch/gpu/}

\textbf{Experimento 6: Comparativa de splits (ImageNet-100)}

\begin{lstlisting}[language=bash]
python3 experiments/imagenet_subset_by_splits/main.py \
    --framework pytorch \
    --processor gpu \
    --model-name vgg16 \
    --initial-layer -2 \
    --epochs 10 \
    --splits 1000 5000 10000
\end{lstlisting}

Resultados en: \path{results/imagenet_subset/gradient_by_splits/pytorch/gpu/}

\textbf{Experimento 7: Aprendizaje continual incremental (ImageNet-100)}

\begin{lstlisting}[language=bash]
python3 experiments/imagenet_continual_learning/main.py \
    --framework pytorch \
    --processor gpu \
    --epochs 10
\end{lstlisting}

Resultados en: \path{results/imagenet_subset/continual_learning/pytorch/gpu/}

\textbf{Experimento 8: Barrido de initial\_layer (ImageNet-100)}

\begin{lstlisting}[language=bash]
python3 experiments/initial_layer_sweep/main.py \
    --framework pytorch \
    --processor gpu
\end{lstlisting}

Resultados en: \path{results/imagenet_subset/initial_layer_sweep/pytorch/gpu/}

\textbf{Experimento 9: Streaming incremental de clases (ImageNet-100)}

\begin{lstlisting}[language=bash]
python3 experiments/imagenet_streaming_incremental/main.py \
    --framework pytorch \
    --processor gpu
\end{lstlisting}

Resultados en: \path{results/imagenet_subset/streaming_incremental/pytorch/gpu/}

\subsection{C.5 Consulta de resultados}

Los resultados de cada ejecución se registran automáticamente en MLflow. Para visualizar el panel de
MLflow:

\begin{lstlisting}[language=bash]
mlflow ui --port 5000
\end{lstlisting}

Acceder en el navegador a \texttt{http://localhost:5000}. Cada experimento aparece bajo su nombre
(\texttt{mnist\_speed\_loss}, \texttt{mnist\_learning\_rate}, etc.) con todos los parámetros de
ejecución registrados.

Las gráficas HTML generadas pueden abrirse directamente en cualquier navegador sin servidor
adicional.

\section{Anexo D: Software desarrollado}

Además de la librería Qsimov (preexistente), este TFG ha requerido el desarrollo de un conjunto de
scripts de experimentación. A continuación se describen los módulos principales desarrollados,
agrupados por experimento.

\subsection{D.1 Módulo de infraestructura de experimentos}

\bpath{experiments/mlflow.py}: Utilidades de orquestación
MLflow: \texttt{run\_script()} (ejecuta subprocesos capturando stdout/stderr y registrando el
resultado en MLflow), \texttt{get\_or\_make\_export\_directory()} (crea el directorio
\texttt{run\_exports/<timestamp>\_<run\_name>/}), \texttt{export\_active\_run()} (serializa el
run activo de MLflow a \texttt{report.json}).

\bpath{experiments/git.py}: Comprueba que el repositorio git esté limpio antes de
lanzar un experimento (\texttt{check\_git\_clean()}), y registra el hash del commit activo en
MLflow para garantizar reproducibilidad.

\bpath{experiments/path_utils.py}: Centraliza las rutas de disco de cada experimento
(directorios de resultados, datasets, directorios de exportación). Cualquier cambio en la
organización de carpetas solo requiere modificar este fichero.

\subsection{D.2 Experimento 1: Velocidad y precisión (MNIST)}

\bpath{experiments/mnist_speed_loss/pytorch/train_pytorch_models.py}: Entrena el
modelo base estándar PyTorch sobre MNIST con distintas fracciones del dataset (\texttt{half},
\texttt{quarter}, \texttt{full}) y dos funciones de pérdida (\texttt{categorical\_crossentropy},
\texttt{mse}). Guarda historia de entrenamiento (accuracy, loss, tiempo por epoch) como
\texttt{.pkl}.

\bpath{experiments/mnist_speed_loss/pytorch/train_qsimov_models.py}: Entrena
\texttt{PytorchQsimovGradient} y \texttt{PytorchQsimovLinearSystem} sobre el mismo dataset.
Comparte el PathSelector del modelo base (mismo \texttt{initial\_layer=-2}).

\bpath{experiments/mnist_speed_loss/speed_loss_plots.py}: Genera gráficas HTML
comparando accuracy y loss en función del tiempo de entrenamiento (eje X = tiempo en segundos) para
los tres métodos. Usa Plotly para generar gráficas interactivas.

\bpath{experiments/mnist_speed_loss/main.py}: Orquestador MLflow: lanza los tres
scripts anteriores en orden, registra parámetros y copia los artefactos generados a
\texttt{run\_exports/}.

\subsection{D.3 Experimento 2: Robustez al learning rate (MNIST)}

\bpath{experiments/mnist_learning_rate/train_pytorch_qsimov_models.py}: Itera
sobre 6 valores de learning rate, entrenando para cada uno tanto el modelo estándar PyTorch como
\texttt{PytorchQsimovGradient}. Registra la accuracy final en test para cada combinación.

\bpath{experiments/mnist_learning_rate/learning_rate_plots.py}: Genera gráficas de
accuracy en función del learning rate (escala logarítmica) y curvas de entrenamiento por epoch para
cada LR.

\bpath{experiments/mnist_learning_rate/main.py}: Orquestador MLflow del experimento.

\subsection{D.4 Experimentos de olvido catastrófico (MNIST y CIFAR-10)}

\bpath{experiments/mnist_forgetting/pytorch_train_base_model.py} y
\bpath{experiments/cifar10_forgetting/pytorch_train_base_model.py}: Entrenan el
modelo base sobre la fase 1 (clases 0--4) y guardan el modelo en disco.

\bpath{experiments/mnist_forgetting/pytorch_run_forgetting.py} y
\bpath{experiments/cifar10_forgetting/pytorch_run_forgetting.py}: Ejecutan todos
los métodos de actualización (QsimovLinearSystem acumulado, QsimovGradient, fine-tuning estándar,
re-entrenamiento acumulado) sobre la fase 2 (clases 5--9) y evalúan la accuracy en ambos conjuntos
de clases.

{\tolerance=9999\relax
\bpath{experiments/mnist_forgetting/plot_forgetting.py} y
\bpath{experiments/cifar10_forgetting/plot_forgetting.py}:
Generan gráficas de barras comparando accuracy en clases antiguas vs.\ nuevas para cada método.
\par}

\subsection{D.5 Experimentos de ImageNet}

\bpath{experiments/imagenet_subset_by_splits/train_pytorch.py}: Entrena el
clasificador estándar VGG16 (base convolucional congelada, top ancho 25088$\rightarrow$4096$\rightarrow$4096$\rightarrow$100)
con distintos tamaños de split.

\bpath{experiments/imagenet_subset_by_splits/train_pytorch_qsimov.py}: Entrena el
PathSelector VGG16 (top estrecho 25088$\rightarrow$2048$\rightarrow$128$\rightarrow$100) y
\texttt{PytorchQsimovGradient} para cada split.

\bpath{experiments/imagenet_continual_learning/train_pytorch_qsimov.py}: Implementa
el escenario de 4 rondas de 25 clases nuevas: entrena \texttt{QsimovLinearSystem} con acumulación
de ecuaciones entre rondas, \texttt{QsimovLinearSystem} con reset, y \texttt{QsimovGradient}.

\bpath{experiments/imagenet_streaming_incremental/train_pytorch_streaming.py}:
Simula el escenario de streaming: 20 batches de 5 clases nuevas cada uno. Para cada batch entrena
los 4 métodos comparados y evalúa sobre el test set acumulado.

\bpath{experiments/initial_layer_sweep/train_pytorch_sweep.py}: Prueba distintos
valores de \texttt{initial\_layer} sobre VGG16, midiendo número de caminos generados, tiempo de
construcción del PathSelector, tiempo de entrenamiento y precisión final.
